% Options for packages loaded elsewhere
\PassOptionsToPackage{unicode}{hyperref}
\PassOptionsToPackage{hyphens}{url}
\PassOptionsToPackage{dvipsnames,svgnames,x11names}{xcolor}
\documentclass[
  american,
  a4paper,
  10pt]{article}
\usepackage{xcolor}
\usepackage[top=22mm,bottom=24mm,left=23mm,right=23mm,headheight=15pt]{geometry}
\usepackage{amsmath,amssymb}
\setcounter{secnumdepth}{-\maxdimen} % remove section numbering
\usepackage{iftex}
\ifPDFTeX
  \usepackage[T1]{fontenc}
  \usepackage[utf8]{inputenc}
  \usepackage{textcomp} % provide euro and other symbols
\else % if luatex or xetex
  \usepackage{unicode-math} % this also loads fontspec
  \defaultfontfeatures{Scale=MatchLowercase}
  \defaultfontfeatures[\rmfamily]{Ligatures=TeX,Scale=1}
\fi
\usepackage{lmodern}
\ifPDFTeX\else
  % xetex/luatex font selection
\fi
% Use upquote if available, for straight quotes in verbatim environments
\IfFileExists{upquote.sty}{\usepackage{upquote}}{}
\IfFileExists{microtype.sty}{% use microtype if available
  \usepackage[]{microtype}
  \UseMicrotypeSet[protrusion]{basicmath} % disable protrusion for tt fonts
}{}
\usepackage{setspace}
\makeatletter
\@ifundefined{KOMAClassName}{% if non-KOMA class
  \IfFileExists{parskip.sty}{%
    \usepackage{parskip}
  }{% else
    \setlength{\parindent}{0pt}
    \setlength{\parskip}{6pt plus 2pt minus 1pt}}
}{% if KOMA class
  \KOMAoptions{parskip=half}}
\makeatother
\usepackage{longtable,booktabs,array}
\newcounter{none} % for unnumbered tables
\usepackage{calc} % for calculating minipage widths
% Correct order of tables after \paragraph or \subparagraph
\usepackage{etoolbox}
\makeatletter
\patchcmd\longtable{\par}{\if@noskipsec\mbox{}\fi\par}{}{}
\makeatother
% Allow footnotes in longtable head/foot
\IfFileExists{footnotehyper.sty}{\usepackage{footnotehyper}}{\usepackage{footnote}}
\makesavenoteenv{longtable}
\ifLuaTeX
\usepackage[bidi=basic,shorthands=off]{babel}
\else
\usepackage[bidi=default,shorthands=off]{babel}
\fi
\ifLuaTeX
  \usepackage{selnolig} % disable illegal ligatures
\fi
\setlength{\emergencystretch}{3em} % prevent overfull lines
\providecommand{\tightlist}{%
  \setlength{\itemsep}{0pt}\setlength{\parskip}{0pt}}
% ARCANA venue-neutral paper style. Keep this file ASCII and engine-neutral.
\usepackage{fancyhdr}
\usepackage{lastpage}
\usepackage{titlesec}
\usepackage{enumitem}
\usepackage{titling}
\usepackage{fvextra}

\definecolor{ArcanaInk}{HTML}{172026}
\definecolor{ArcanaGreen}{HTML}{1F6F5B}
\definecolor{ArcanaRule}{HTML}{C9D2D0}
\definecolor{ArcanaPale}{HTML}{F2F6F5}

\color{ArcanaInk}
\setlength{\headheight}{15pt}
\setlength{\droptitle}{-1.4cm}
\setlength{\parindent}{0pt}
\setlength{\parskip}{5pt plus 1pt minus 1pt}
\setlength{\LTpre}{5pt}
\setlength{\LTpost}{7pt}
\renewcommand{\arraystretch}{1.12}
\raggedbottom
\widowpenalty=10000
\clubpenalty=10000
\displaywidowpenalty=10000

\pagestyle{fancy}
\fancyhf{}
\fancyhead[L]{\small\sffamily ARCANA v1.0}
\fancyhead[R]{\small\sffamily Open research/reference}
\fancyfoot[C]{\small\sffamily \thepage\ / \pageref*{LastPage}}
\renewcommand{\headrulewidth}{0.4pt}
\renewcommand{\footrulewidth}{0pt}

\fancypagestyle{plain}{
  \fancyhf{}
  \fancyfoot[C]{\small\sffamily \thepage\ / \pageref*{LastPage}}
  \renewcommand{\headrulewidth}{0pt}
}

\titleformat{\section}
  {\Large\bfseries\sffamily\color{ArcanaInk}}
  {}{0pt}{}
  [\vspace{2pt}\color{ArcanaRule}\titlerule]
\titleformat{\subsection}
  {\large\bfseries\sffamily\color{ArcanaGreen}}
  {}{0pt}{}
\titleformat{\subsubsection}
  {\normalsize\bfseries\sffamily\color{ArcanaInk}}
  {}{0pt}{}
\titlespacing*{\section}{0pt}{18pt}{8pt}
\titlespacing*{\subsection}{0pt}{13pt}{5pt}
\titlespacing*{\subsubsection}{0pt}{10pt}{3pt}

\pretitle{\begin{flushleft}\color{ArcanaInk}\Huge\bfseries\sffamily}
\posttitle{\par\vspace{6pt}\color{ArcanaGreen}\rule{\textwidth}{1.2pt}\end{flushleft}}
\preauthor{\begin{flushleft}\large\sffamily}
\postauthor{%
  \par\vspace{2pt}%
  {\small TradePhantom LLC\\
  \href{mailto:dev@tradephantom.com}{\texttt{dev@tradephantom.com}}\\
  \href{https://getaxcp.com}{\texttt{getaxcp.com}}}%
  \par\end{flushleft}%
}
\predate{\begin{flushleft}\small\sffamily\color{ArcanaGreen}}
\postdate{%
  \par\vspace{3pt}%
  \href{https://doi.org/10.5281/zenodo.21333463}{%
    \path{doi:10.5281/zenodo.21333463}}%
  \par\end{flushleft}\vspace{4pt}%
}

\setlist[itemize]{leftmargin=1.35em,itemsep=2pt,topsep=3pt}
\setlist[enumerate]{leftmargin=1.55em,itemsep=2pt,topsep=3pt}
\setcounter{tocdepth}{2}

\AtBeginEnvironment{longtable}{\small}
\AtBeginDocument{%
  \fvset{breaklines=true,breakanywhere=true,fontsize=\small}%
  \hypersetup{pdfinfo={%
    DOI={10.5281/zenodo.21333463},%
    License={CC BY 4.0},%
    Website={https://getaxcp.com}%
  }}%
}
\usepackage{bookmark}
\IfFileExists{xurl.sty}{\usepackage{xurl}}{} % add URL line breaks if available
\urlstyle{same}
% fallback for those not using the hyperref driver hyperxmp:
\makeatletter
\@ifundefined{xmpquote}{\newcommand{\xmpquote}[1]{#1}}{}
\makeatother
\hypersetup{
  pdftitle={ARCANA: Autonomy Risk Calculus for Agentic Network Assurance},
  pdfauthor={Julio Elizondo Rodriguez},
  pdflang={en-US},
  pdfsubject={Open research preprint on model-bounded autonomy accounting for agentic
systems},
  pdfkeywords={\xmpquote{agentic systems}, \xmpquote{autonomy
accounting}, \xmpquote{model-bounded risk}, \xmpquote{capability
graphs}, \xmpquote{calibration uncertainty}, \xmpquote{FastGate}},
  colorlinks=true,
  linkcolor={ArcanaGreen},
  filecolor={Maroon},
  citecolor={ArcanaGreen},
  urlcolor={ArcanaGreen},
  pdfcreator={LaTeX via pandoc}}

\ifXeTeX
\special{pdf:trailerid [ <7c07dbb952113d4a604dbaecf2db8afd> <7c07dbb952113d4a604dbaecf2db8afd> ]}
\fi
\ifPDFTeX
\pdftrailerid{}
\pdftrailer{/ID [ <7c07dbb952113d4a604dbaecf2db8afd> <7c07dbb952113d4a604dbaecf2db8afd> ]}
\fi
\ifLuaTeX
\pdfvariable trailerid {[ <7c07dbb952113d4a604dbaecf2db8afd> <7c07dbb952113d4a604dbaecf2db8afd> ]}
\fi

\title{ARCANA: Autonomy Risk Calculus for Agentic Network Assurance}
\usepackage{etoolbox}
\makeatletter
\providecommand{\subtitle}[1]{% add subtitle to \maketitle
  \apptocmd{\@title}{\par {\large #1 \par}}{}{}
}
\makeatother
\subtitle{Model-Bounded Autonomy Accounting for Agentic Systems}
\author{Julio Elizondo Rodriguez}
\date{13 July 2026}

\begin{document}
\maketitle

\noindent
\colorbox{ArcanaPale}{%
  \parbox{\dimexpr\textwidth-2\fboxsep\relax}{%
    \small\sffamily
    \textbf{Public research manuscript v1.0, post-hardening repository edition.}
    Reference implementation and deterministic synthetic evaluator evidence
    only. This artifact is not a production approval, enforcement approval,
    external peer-review claim, or commercial certificate authorization. The
    DOI identifies this preprint; it does not elevate its evidence class.
  }%
}
\vspace{10pt}

\setstretch{1.05}
\begin{abstract}
Agentic systems can act through tools, memory, delegation, capability envelopes,
repeated operations, and external services. Task success alone does not show whether a
proposed operation remains inside an acceptable autonomy-risk envelope. ARCANA, the
Autonomy Risk Calculus for Agentic Network Assurance, introduces a public
research/reference framework for autonomy accounting: it estimates whether a proposed
operation remains bounded under an explicit risk model version, calibration profile,
decision horizon, uncertainty bounds, evidence assumptions, graph state, and reason-code
contract.

ARCANA models unsafe propagation pressure separately from loss or impact. It represents
an agentic system as a capability graph, assigns horizon-bound propagation weights,
computes upper-bound spectral propagation risk, evaluates loss-side quantities
separately, and emits schema-shaped risk contexts or certificate-like artifacts with
explicit caveats. The public reference implementation uses deterministic synthetic
fixtures, public schemas, public reason codes, and local reproducibility gates.

This manuscript does not claim production readiness, production enforcement, universal
safety, or production certificate issuance. A0 outputs remain non-certifiable. Benchmark
scenarios are synthetic and are intended to show why task success, unsafe action rate,
risk-bound status, and artifact-contract validity must be reported separately. ARCANA is
only as complete as the graph and evidence it is given.
\end{abstract}

\clearpage
\tableofcontents
\clearpage

\section{1. Paper Rule}\label{1-paper-rule}

ARCANA may be described as:

\begin{Verbatim}[breaklines=true,breakanywhere=true,fontsize=\small]
model-bounded autonomy accounting for agentic systems
\end{Verbatim}

The strongest allowed public claim in this manuscript is:

\begin{Verbatim}[breaklines=true,breakanywhere=true,fontsize=\small]
ARCANA estimates bounded autonomy under explicit risk model versions, calibration profiles, decision horizons, uncertainty bounds, and evidence assumptions.
\end{Verbatim}

The manuscript must not present any agent, workflow, deployment, model, policy, or
organization as safe by proof. Every result is scoped by declared assumptions, public
schemas, reason codes, evidence, and limitations.

\section{Contributions}\label{contributions}

ARCANA contributes:

\begin{enumerate}
\def\labelenumi{\arabic{enumi}.}
\tightlist
\item
  A graph-based model for autonomy-risk propagation in agentic systems.
\item
  A strict separation between propagation risk \texttt{K} and loss or impact \texttt{L}.
\item
  A decision-horizon contract that prevents dimensionally ambiguous risk claims.
\item
  A calibration-level model that exposes uncertainty and evidence quality.
\item
  A conservative upper-bound spectral decision rule.
\item
  FastGate, a fail-closed sparse-update path for low-latency admission.
\item
  Certificate-like bounded artifacts with reason-code semantics.
\item
  ARCANA-Bench, a synthetic benchmark separating task success from autonomy risk.
\end{enumerate}

\section{2. Problem}\label{2-problem}

Agentic systems can complete tasks while increasing autonomy risk. Examples include:

\begin{itemize}
\tightlist
\item
  routing unsafe instructions through an otherwise useful tool path;
\item
  retaining poisoned memory that affects future operations;
\item
  requesting a tool scope broader than the task requires;
\item
  delegating across agents until risk compounds;
\item
  optimizing for benchmark success while hiding unsafe actions;
\item
  generating dynamic behavior outside the calibrated envelope.
\end{itemize}

These behaviors are not captured by a task-success metric alone. ARCANA asks a narrower
question:

\begin{Verbatim}[breaklines=true,breakanywhere=true,fontsize=\small]
Given a proposed operation, graph state, calibration profile, decision horizon, uncertainty interval, and evidence bundle, does the operation remain inside the declared bounded-risk envelope?
\end{Verbatim}

This is not the same question as whether a local policy should allow the operation.
ARCANA can inform surrounding systems, but the surrounding system remains responsible
for local policy, enforcement, and operational approval.

\section{3. Scope and Non-Scope}\label{3-scope-and-non-scope}

Public ARCANA includes:

\begin{itemize}
\tightlist
\item
  formal model documentation;
\item
  calibration methodology;
\item
  reason-code registry;
\item
  public schemas;
\item
  reference implementation;
\item
  synthetic demos;
\item
  ARCANA-Bench synthetic scenarios;
\item
  public integration contract;
\item
  limitations and release-readiness documents.
\end{itemize}

Public ARCANA does not include:

\begin{itemize}
\tightlist
\item
  production enforcement implementation;
\item
  private policy engine;
\item
  customer telemetry pipeline;
\item
  production dashboard;
\item
  private deployment profile;
\item
  commercial certificate issuance;
\item
  private evidence payload;
\item
  production operational approval.
\end{itemize}

The public repository is a research/reference track. It is not a production control
plane.

ARCANA is protocol-independent and can be used without AXCP. Within the related AXCP
project family, AXCP addresses secure message exchange while ARCANA addresses
model-bounded autonomy accounting. The two projects remain separable: ARCANA is not an
AXCP Core dependency, and this paper does not make AXCP a precondition for the ARCANA
model or reference implementation {[}1{]}.

\section{4. Model Overview}\label{4-model-overview}

ARCANA represents an agentic system at time \texttt{t} as a graph:

\begin{Verbatim}[breaklines=true,breakanywhere=true,fontsize=\small]
G_t = (V_t, E_t)
\end{Verbatim}

where \path{V_t} contains public node classes:

\begin{Verbatim}[breaklines=true,breakanywhere=true,fontsize=\small]
A  = agents
T  = tools or APIs
M  = memory or knowledge stores
H  = human or policy gates
Gt = grants or capability envelopes
O  = operations or plans
D  = distilled or repeated operations
X  = external systems
\end{Verbatim}

An edge is:

\begin{Verbatim}[breaklines=true,breakanywhere=true,fontsize=\small]
e = (source, target, edge_type, factors, evidence)
\end{Verbatim}

Each edge describes a possible unsafe propagation path over a declared decision horizon.
The graph hash binds matrix ordering and artifact scope to the graph state used for
evaluation.

\section{5. Decision Horizon}\label{5-decision-horizon}

A decision horizon is:

\begin{Verbatim}[breaklines=true,breakanywhere=true,fontsize=\small]
H = (horizon_id, duration_seconds, context)
\end{Verbatim}

Every model input must be horizon-compatible:

\begin{itemize}
\tightlist
\item
  graph edge weights;
\item
  calibration profile;
\item
  risk context;
\item
  certificate-like artifact;
\item
  proposed graph delta;
\item
  loss bounds where loss is in scope.
\end{itemize}

Without \texttt{H}, propagation weights are dimensionally ambiguous. Missing or
incompatible horizons return:

\begin{Verbatim}[breaklines=true,breakanywhere=true,fontsize=\small]
ARCANA_DENY_DECISION_HORIZON_MISMATCH
\end{Verbatim}

The horizon is not a formatting detail. It is part of the mathematical contract.

\section{6. Edge Weights and Propagation
Matrix}\label{6-edge-weights-and-propagation-matrix}

For an edge \texttt{e} over horizon \texttt{H}, ARCANA models unsafe propagation weight
as:

\begin{Verbatim}[breaklines=true,breakanywhere=true,fontsize=\small]
w_e(H) =
  activation_e(H)
  * p_unsafe_e(H)
  * exposure_e(H)
  * capability_e(H)
  * scope_e(H)
  * (1 - detectability_e(H))
  * (1 - gate_effectiveness_e(H))
  * (1 - reversibility_e(H))
\end{Verbatim}

Risk-amplifying factors use conservative upper bounds when unknown. Control factors use
conservative lower bounds when unknown.

Expected factor domains:

{\def\LTcaptype{none} % do not increment counter
\begin{longtable}[]{@{}>{\raggedright\arraybackslash}p{0.27\linewidth}>{\raggedright\arraybackslash}p{0.17\linewidth}>{\raggedright\arraybackslash}p{0.47\linewidth}@{}}
\toprule\noalign{}
Factor & Domain & Meaning \\
\midrule\noalign{}
\endhead
\bottomrule\noalign{}
\endlastfoot
\path{activation_e(H)} & \texttt{\textgreater{}=\ 0} & Expected unsafe transition
pressure over horizon \texttt{H}; may exceed \texttt{1} only when repeated transition
pressure is intentionally modeled. \\
\path{p_unsafe_e(H)} & \texttt{{[}0,\ 1{]}} & Conditional unsafe-transition
likelihood under the declared model. \\
\path{exposure_e(H)} & \texttt{{[}0,\ 1{]}} & Degree to which the target can be
reached by the source under the current operational context. \\
\path{capability_e(H)} & \texttt{{[}0,\ 1{]}} & Privilege amplification or capability
strength if the target is reached. \\
\path{scope_e(H)} & \texttt{{[}0,\ 1{]}} & Breadth of resources or actions available
through the edge. \\
\path{detectability_e(H)} & \texttt{{[}0,\ 1{]}} & Strength of detection over the
horizon. \\
\path{gate_effectiveness_e(H)} & \texttt{{[}0,\ 1{]}} & Strength of human, policy,
or technical gate controls. \\
\path{reversibility_e(H)} & \texttt{{[}0,\ 1{]}} & Degree to which the transition can
be undone or contained. \\
\end{longtable}
}

Calibration profiles must prevent double-counting between \path{exposure_e},
\path{capability_e}, and \path{scope_e}. Each factor should have a definition,
domain, source, uncertainty rule, conservative default direction, and caveat when
evidence is sparse or stale.

For \texttt{n\ =\ \textbar{}V\_t\textbar{}}, the propagation matrix is:

\begin{Verbatim}[breaklines=true,breakanywhere=true,fontsize=\small]
K(H) in R_nonnegative^(n x n)
\end{Verbatim}

with:

\begin{Verbatim}[breaklines=true,breakanywhere=true,fontsize=\small]
K_ij(H) = sum w_e(H) for all edges e where source(e) = i and target(e) = j
\end{Verbatim}

\texttt{K} must be square, nonnegative, horizon-bound, and graph-bound. It contains
unsafe propagation pressure only. If graph, matrix, interval, threshold, node ordering,
or required model input is malformed, ARCANA returns:

\begin{Verbatim}[breaklines=true,breakanywhere=true,fontsize=\small]
ARCANA_DENY_MODEL_INPUT_INVALID
\end{Verbatim}

\section{7. Uncertainty and Spectral Risk}\label{7-uncertainty-and-spectral-risk}

ARCANA tracks matrix uncertainty as:

\begin{Verbatim}[breaklines=true,breakanywhere=true,fontsize=\small]
K_lower(H), K_mean(H), K_upper(H)
\end{Verbatim}

with the invariant:

\begin{Verbatim}[breaklines=true,breakanywhere=true,fontsize=\small]
0 <= K_lower_ij <= K_mean_ij <= K_upper_ij
\end{Verbatim}

ARCANA computes:

\begin{Verbatim}[breaklines=true,breakanywhere=true,fontsize=\small]
rho_lower = rho(K_lower)
rho_mean  = rho(K_mean)
rho_upper = rho(K_upper)
\end{Verbatim}

Admission-like decisions use \path{rho_upper}, not \path{rho_mean}.

The public subcritical-under-model condition is:

\begin{Verbatim}[breaklines=true,breakanywhere=true,fontsize=\small]
0 <= theta_rho <= 1
rho(K_upper) < theta_rho
\end{Verbatim}

where \path{theta_rho} is the declared public threshold for the decision context. A
stricter policy threshold may be used only inside the same subcritical range:

\begin{Verbatim}[breaklines=true,breakanywhere=true,fontsize=\small]
rho(K_upper) < theta_rho <= 1
\end{Verbatim}

This condition means the declared model is subcritical under the stated assumptions. A
threshold above \texttt{1} may be a diagnostic or policy quantity in some other model,
but it must not support an ARCANA public allow-like, subcritical-under-model, or
certificate-like bounded-autonomy artifact. Public allow-like, subcritical-under-model,
and certificate-like bounded-autonomy artifacts must not use
\texttt{theta\_rho\ \textgreater{}\ 1}. It is not a statement about all possible
real-world behavior.

When the upper-bound propagation risk exceeds the threshold or budget, ARCANA returns:

\begin{Verbatim}[breaklines=true,breakanywhere=true,fontsize=\small]
ARCANA_DENY_RHO_UPPER_BOUND
\end{Verbatim}

\section{Mathematical Propositions and Proof
Sketches}\label{mathematical-propositions-and-proof-sketches}

The public model relies on standard facts about nonnegative matrices. These propositions
are model-bounded statements, not claims about all real-world behavior.

The nonnegative-matrix results used below are standard consequences of Perron-Frobenius
and Collatz-Wielandt theory {[}5, 8{]}. The subcritical propagation interpretation is an
analogy to finite-type branching-process models {[}6{]}, not a claim that an agentic
system satisfies every independence assumption of a classical branching process.

\subsection{Proposition 1 - Subcritical finite
propagation}\label{proposition-1---subcritical-finite-propagation}

For a nonnegative propagation matrix \texttt{K} with:

\begin{Verbatim}[breaklines=true,breakanywhere=true,fontsize=\small]
rho(K) < 1
\end{Verbatim}

the Neumann series converges:

\begin{Verbatim}[breaklines=true,breakanywhere=true,fontsize=\small]
sum_{m=0}^{infinity} K^m
\end{Verbatim}

Under the declared ARCANA model, this means modeled propagation pressure remains finite.
It does not prove that unmodeled propagation paths are absent.

\subsection{Proposition 2 - Upper-bound
monotonicity}\label{proposition-2---upper-bound-monotonicity}

If:

\begin{Verbatim}[breaklines=true,breakanywhere=true,fontsize=\small]
0 <= K <= K_upper
\end{Verbatim}

entrywise, then:

\begin{Verbatim}[breaklines=true,breakanywhere=true,fontsize=\small]
rho(K) <= rho(K_upper)
\end{Verbatim}

This justifies using \path{rho_upper} for admission-like decisions when uncertainty
is represented through entrywise upper-bound matrices.

\subsection{Proposition 3 - FastGate conservative
bound}\label{proposition-3---fastgate-conservative-bound}

For \texttt{A\ \textgreater{}=\ 0} and \texttt{x\ \textgreater{}\ 0}, the
Collatz-Wielandt upper bound gives:

\begin{Verbatim}[breaklines=true,breakanywhere=true,fontsize=\small]
rho(A) <= max_i ((A x)_i / x_i)
\end{Verbatim}

FastGate may use this only as a conservative upper bound. If the vector, nonnegativity,
graph, cache, margin, calibration, or horizon assumptions fail, FastGate must not return
an allow-like result.

\section{8. Loss Is Separate From
Propagation}\label{8-loss-is-separate-from-propagation}

ARCANA separates:

\begin{Verbatim}[breaklines=true,breakanywhere=true,fontsize=\small]
K = unsafe propagation pressure
L = impact or loss
\end{Verbatim}

Loss can include operational, financial, privacy, compliance, safety, or reputational
impact. Loss must not be embedded in \texttt{K}; propagation risk must not be treated as
impact.

When loss is in scope, ARCANA can report loss-side upper bounds such as:

\begin{Verbatim}[breaklines=true,breakanywhere=true,fontsize=\small]
AaR_99_upper
AES_99_upper
\end{Verbatim}

If a required loss model is missing, malformed, unsupported, or horizon-incompatible,
ARCANA returns:

\begin{Verbatim}[breaklines=true,breakanywhere=true,fontsize=\small]
ARCANA_DENY_LOSS_MODEL_INVALID
\end{Verbatim}

If the upper-bound Autonomy-at-Risk or Agentic Expected Shortfall exceeds declared
limits, ARCANA returns:

\begin{Verbatim}[breaklines=true,breakanywhere=true,fontsize=\small]
ARCANA_DENY_AAR_UPPER_BOUND
ARCANA_DENY_AES_UPPER_BOUND
\end{Verbatim}

An operation can have bounded propagation risk but unacceptable loss bounds. It can also
have low expected impact but unacceptable propagation-risk upper bound.

The tail-loss terminology is informed by the Value-at-Risk and conditional Value-at-Risk
literature {[}7{]}. ARCANA does not import a financial loss model into \texttt{K}; it
applies loss-side quantities only through the separately declared \texttt{L} contract.

\section{9. Calibration}\label{9-calibration}

Calibration converts evidence and assumptions into uncertainty bounds. It is not a proof
of real-world behavior.

ARCANA uses four public calibration levels:

{\def\LTcaptype{none} % do not increment counter
\begin{longtable}[]{@{}>{\raggedright\arraybackslash}p{0.07\linewidth}>{\raggedright\arraybackslash}p{0.19\linewidth}>{\raggedright\arraybackslash}p{0.34\linewidth}>{\raggedright\arraybackslash}p{0.27\linewidth}@{}}
\toprule\noalign{}
Level & Name & Public use & Certification status \\
\midrule\noalign{}
\endhead
\bottomrule\noalign{}
\endlastfoot
A0 & Heuristic / Demo Only & Synthetic examples, tutorials, paper illustrations. &
Always \path{non_certifiable}. \\
A1 & Static Conservative Prior & Conservative offline estimates from declared static
factors. & Public v0.2 default is \path{non_certifiable}; no certificate-like
artifact support. \\
A2 & Empirical / Red-Team Calibration & Controlled adversarial evidence, reviewed
replay, or empirical public benchmark evidence. & Eligible for
\path{certifiable_under_profile} only after reviewed scope, coverage, freshness, and
non-synthetic evidence checks. \\
A3 & Runtime Bayesian Calibration & Maintained segmented distributions with decay and
priors. & Eligible for \path{certifiable_under_profile} only after reviewed runtime
evidence and auditability checks. \\
\end{longtable}
}

A0 is demo-only. A0 outputs must preserve:

\begin{Verbatim}[breaklines=true,breakanywhere=true,fontsize=\small]
ARCANA_INFO_A0_NON_CERTIFIABLE
\end{Verbatim}

The public claim boundary after production-readiness review is:

\begin{itemize}
\tightlist
\item
  A1 is a conservative prior level and is not a commercial certificate basis.
\item
  A1 can support bounded reference risk contexts when all other constraints pass, but A1
  remains \path{non_certifiable} and cannot support non-demo certificate-like
  artifacts in the public v0.2 contract.
\item
  A2 may support reviewed empirical calibration claims only when scope, evidence,
  freshness, horizon, and non-synthetic evidence match the reviewed profile.
\item
  A3 remains a runtime-calibration design target unless a separate runtime distribution
  review is completed.
\item
  Public examples do not disclose private thresholds, formulas, runtime distributions,
  telemetry, or deployment mechanics.
\item
  No public calibration level in this manuscript authorizes production enforcement or
  certificate issuance by itself.
\end{itemize}

The public source classes deliberately distinguish \path{public_benchmark} from
\path{empirical_public_benchmark}. The first names the current synthetic
ARCANA-Bench reference-evaluator evidence and never qualifies A2 or a certifiable state.
The second is reserved for non-synthetic public benchmark evidence with separately
reviewed provenance, scope, coverage, freshness, and reproducibility. Relabeling a
synthetic run cannot upgrade its calibration level.

The phrase \path{certifiable_under_profile} is a bounded public/reference semantic
state. It means that an artifact may support a bounded-under-profile claim inside the
declared ARCANA model and calibration profile. It does not imply commercial certificate
issuance, production enforcement, customer approval, or operational authorization.

If the requested decision requires stronger calibration than the supplied profile,
ARCANA returns:

\begin{Verbatim}[breaklines=true,breakanywhere=true,fontsize=\small]
ARCANA_DENY_CALIBRATION_INSUFFICIENT
\end{Verbatim}

If evaluation is allowed but admission-like claims are not, ARCANA may return:

\begin{Verbatim}[breaklines=true,breakanywhere=true,fontsize=\small]
ARCANA_REQUIRE_OBSERVE_ONLY
\end{Verbatim}

Calibration limitations must be visible in the result. Sparse evidence widens
uncertainty. Stale evidence weakens or invalidates the profile. Passive absence of
observed incidents cannot collapse uncertainty by itself.

\section{10. Evidence Binding}\label{10-evidence-binding}

ARCANA artifacts bind results to evidence by source identifiers and hashes. A public
evidence reference may include:

\begin{Verbatim}[breaklines=true,breakanywhere=true,fontsize=\small]
source_type
source_id
evidence_hash
synthetic
\end{Verbatim}

For public examples, \texttt{synthetic} must be true. An evidence hash binds the
artifact to evidence; it does not prove evidence quality by itself.

Evidence mismatch, stale evidence, missing evidence, or unsupported evidence class can
invalidate the result. The correct action is recompute, denial, or observe-only
according to the public reason-code contract.

\section{11. Decision Semantics}\label{11-decision-semantics}

ARCANA emits explicit verdicts:

{\def\LTcaptype{none} % do not increment counter
\begin{longtable}[]{@{}>{\raggedright\arraybackslash}p{0.31\linewidth}>{\raggedright\arraybackslash}p{0.62\linewidth}@{}}
\toprule\noalign{}
Verdict & Meaning \\
\midrule\noalign{}
\endhead
\bottomrule\noalign{}
\endlastfoot
\path{allow_bounded_autonomy} & The proposed operation is bounded under declared
assumptions. \\
\path{allow_with_controls} & The operation is bounded only if listed controls are
applied. \\
\path{require_scope_reduction} & The requested scope is too broad but may become
bounded if narrowed. \\
\path{require_human_gate} & A human or policy gate is required before admission. \\
\path{observe_only} & The result may be analyzed but not treated as bounded for
admission or certification. \\
\texttt{deny} & The operation is outside the declared bounded-risk envelope or
invalid. \\
\end{longtable}
}

Allow-like results require compatible model, calibration, horizon, graph, evidence,
uncertainty, threshold, budget, and controls.

Distinct failures must not collapse into generic denial. Examples:

\begin{Verbatim}[breaklines=true,breakanywhere=true,fontsize=\small]
ARCANA_DENY_RISK_MODEL_UNSUPPORTED
ARCANA_DENY_CONTEXT_STALE
ARCANA_DENY_GRAPH_HASH_MISMATCH
ARCANA_DENY_BUDGET_EXHAUSTED
ARCANA_REQUIRE_SCOPE_REDUCTION
ARCANA_REQUIRE_HUMAN_GATE
\end{Verbatim}

Reason codes are part of the public contract surface and must be preserved by
integrations.

\section{12. Certificate-Like Artifacts}\label{12-certificate-like-artifacts}

Certificate-like artifacts are bounded records. They are not broad approvals.

Every certificate-like artifact must include:

\begin{itemize}
\tightlist
\item
  schema version;
\item
  certificate identifier;
\item
  certificate type;
\item
  risk model version;
\item
  calibration profile;
\item
  calibration level;
\item
  decision horizon;
\item
  graph hash;
\item
  uncertainty bounds;
\item
  evidence source or hash;
\item
  verdict;
\item
  reason codes;
\item
  caveats;
\item
  issued time and expiry when applicable.
\end{itemize}

A0 artifacts are non-certifiable and cannot support production certificate issuance.

In this public manuscript, A1 certificate-like artifacts are not supported beyond
non-certifiable reference contexts. A2 and A3 references are conditional claim
boundaries: they describe what a reviewed profile would need to bind, not an
authorization to issue a production certificate from this repository. Certificate
issuance, if ever implemented by a deployment-specific system, requires a separate
issuance process, expiry, revocation path, support evidence, and review record.

If a certificate-like artifact is stale, graph-mismatched, horizon-mismatched,
evidence-mismatched, or missing required fields, it must be rejected or recomputed.

\section{13. FastGate}\label{13-fastgate}

FastGate is a conservative optimization path for sparse graph changes. It must preserve
the exact upper-bound decision rule:

\begin{Verbatim}[breaklines=true,breakanywhere=true,fontsize=\small]
0 <= theta_rho <= 1
rho(K_after_upper) < theta_rho
\end{Verbatim}

FastGate may return an allow-like result only when it proves a conservative bound:

\begin{Verbatim}[breaklines=true,breakanywhere=true,fontsize=\small]
fastgate_upper_bound >= rho(K_after_upper)
fastgate_upper_bound < theta_rho
\end{Verbatim}

If FastGate cannot prove both inequalities, it must not return
\path{allow_bounded_autonomy}.

Recognized public modes:

{\def\LTcaptype{none} % do not increment counter
\begin{longtable}[]{@{}>{\raggedright\arraybackslash}p{0.27\linewidth}>{\raggedright\arraybackslash}p{0.17\linewidth}>{\raggedright\arraybackslash}p{0.47\linewidth}@{}}
\toprule\noalign{}
Mode & Purpose & Allow-like result permitted \\
\midrule\noalign{}
\endhead
\bottomrule\noalign{}
\endlastfoot
\path{exact_recompute} & Full recompute of \path{rho(K_after_upper)}. & Yes, if
exact constraints pass. \\
\path{perron_collatz_bound} & Conservative sparse update bound. & Yes, if bound
proves admission with margin. \\
\path{observe_only} & Evaluate or explain without admission authority. & No. \\
\end{longtable}
}

FastGate must validate positive-vector assumptions, reducible graph handling, cache
keys, numerical margin, and evidence binding.

The public FastGate implementation also recomputes a domain-separated canonical hash
over each sparse delta and requires a current evidence hash for every admission-capable
path. Its Collatz products, sums, and ratios are rounded upward before the separate
numerical margin is applied. These controls are reference semantics, not a production
latency or enforcement claim.

If the fast path is inconclusive:

\begin{Verbatim}[breaklines=true,breakanywhere=true,fontsize=\small]
ARCANA_DENY_FASTGATE_UNCERTAIN
\end{Verbatim}

If positive-vector requirements are invalid:

\begin{Verbatim}[breaklines=true,breakanywhere=true,fontsize=\small]
ARCANA_DENY_FASTGATE_VECTOR_INVALID
\end{Verbatim}

FastGate is not a weaker model. It is a faster path that must fail closed when
assumptions fail.

\section{14. Distillation Risk}\label{14-distillation-risk}

Distillation means treating an operation, plan, or behavior as repeatable. This can
compress repeated autonomy into a persistent pattern.

Distillation is not eligible unless evidence, stability, scope, calibration, and
invalidation rules are sufficient. If they are not sufficient, ARCANA returns:

\begin{Verbatim}[breaklines=true,breakanywhere=true,fontsize=\small]
ARCANA_DENY_DISTILLATION_RISK
\end{Verbatim}

Public examples must not imply that a demo behavior can become a production reusable
pattern without reviewed calibration evidence.

\section{15. ARCANA-Bench}\label{15-arcana-bench}

ARCANA-Bench v0.1 demonstrates why autonomy risk accounting differs from task-success
evaluation. The public benchmark suite uses synthetic scenarios only.

Scenario classes:

{\def\LTcaptype{none} % do not increment counter
\begin{longtable}[]{@{}>{\raggedright\arraybackslash}p{0.31\linewidth}>{\raggedright\arraybackslash}p{0.62\linewidth}@{}}
\toprule\noalign{}
Scenario class & Benchmark question \\
\midrule\noalign{}
\endhead
\bottomrule\noalign{}
\endlastfoot
Indirect prompt injection containment & Can task success coexist with unsafe
instruction-routing risk? \\
Persistent memory poisoning containment & Can stale or poisoned memory affect future
autonomy risk? \\
Tool misuse scope reduction & Does requested tool scope exceed the bounded envelope? \\
Unsafe delegation cascade & Does delegation amplify propagation risk across agents? \\
Benchmark gaming detection & Can high task success hide unsafe action behavior? \\
Dynamic execution risk & Does runtime-generated behavior exceed calibration scope? \\
\end{longtable}
}

Required metrics:

\begin{Verbatim}[breaklines=true,breakanywhere=true,fontsize=\small]
task_success_rate
unsafe_action_rate
policy_violation_rate
delta_rho_upper
rho_peak
AaR_99_upper
AES_99_upper
containment_time_steps
autonomy_budget_consumed
human_intervention_efficiency
artifact_contract_validity_rate
\end{Verbatim}

Task success is not an ARCANA safety score. Benchmark reports must separate at least:

\begin{Verbatim}[breaklines=true,breakanywhere=true,fontsize=\small]
task_success_rate
unsafe_action_rate
risk_bound_status
artifact_contract_validity_rate
\end{Verbatim}

\path{artifact_contract_validity_rate} measures whether benchmark scenario and
execution artifacts satisfy required public schema and typed contracts. The v0.1 runner
does not emit or validate a production certificate. This metric does not mean production
certification or operational approval.

Earlier drafts used the name \path{certificate_validity_rate}; this manuscript uses
\path{artifact_contract_validity_rate} to avoid implying that the metric validates
a production certificate.

\subsection{Negative Controls and Failure
Fixtures}\label{negative-controls-and-failure-fixtures}

ARCANA-Bench v0.1 includes public synthetic negative-control fixtures that declare
expected fail-closed outcomes for synthetic failure modes. The current runner validates
fixture and execution schemas, constructs typed inputs, invokes the actual public
reference evaluator, and compares observed verdicts, ordered reason codes, and required
controls with the declared contracts. This verifies deterministic reference-evaluator
behavior for those inputs; it does not execute an agentic workload or establish
empirical fail-closed performance. Each negative-control fixture declares
\texttt{scenario\_type:\ negative\_control} and a required \path{negative_control}
identifier.

The runner invokes the actual public reference evaluator; it does not copy the expected
fixture outcome into the observed report.

{\def\LTcaptype{none} % do not increment counter
\begin{longtable}[]{@{}>{\raggedright\arraybackslash}p{0.31\linewidth}>{\raggedright\arraybackslash}p{0.62\linewidth}@{}}
\toprule\noalign{}
Negative-control identifier & Required behavior \\
\midrule\noalign{}
\endhead
\bottomrule\noalign{}
\endlastfoot
\path{missing_decision_horizon} & deny with
\path{ARCANA_DENY_DECISION_HORIZON_MISMATCH}. \\
\path{graph_hash_mismatch} & deny with
\path{ARCANA_DENY_GRAPH_HASH_MISMATCH}. \\
\path{a0_artifact_used_for_admission} & observe-only or deny; never production
admission. \\
\path{rho_mean_below_threshold_rho_upper_above_threshold} & must not allow;
current public fixture denies with \path{ARCANA_DENY_RHO_UPPER_BOUND}. \\
\path{fastgate_inconclusive} & fail closed with
\path{ARCANA_DENY_FASTGATE_UNCERTAIN}. \\
\path{loss_model_missing_required} & deny with
\path{ARCANA_DENY_LOSS_MODEL_INVALID}. \\
\path{stale_evidence} & recompute, observe-only, or deny according to reason-code
contract; current public fixture denies with \path{ARCANA_DENY_CONTEXT_STALE}. \\
\end{longtable}
}

\section{16. Reference Implementation}\label{16-reference-implementation}

The public reference implementation is a local, deterministic implementation of the
public contracts. It is intended for review, reproduction, and tests, not production
deployment.

Public implementation components include:

\begin{itemize}
\tightlist
\item
  typed graph and calibration objects;
\item
  schema validation;
\item
  spectral risk calculator;
\item
  decision evaluator with explicit reason-code branches;
\item
  A0 certificate-like demo output marked non-certifiable;
\item
  FastGate prototype;
\item
  canonical sparse-delta hashing and evidence-bound FastGate admission;
\item
  ARCANA-Bench synthetic fixture loader;
\item
  ARCANA-Bench negative-control coverage validator;
\item
  ARCANA-Bench typed execution-suite loader and actual evaluator runner;
\item
  deterministic machine-readable report with metric provenance and content hashes;
\item
  public audit and validation tools.
\end{itemize}

Local reproduction:

\begin{Verbatim}[breaklines=true,breakanywhere=true,fontsize=\small]
python3 -m venv .venv
. .venv/bin/activate
python -m pip install -e ".[dev]"
make check
make test
make demo
make bench
\end{Verbatim}

\texttt{make\ check} includes public-boundary audit, schema/example validation, and the
deterministic synthetic evaluator benchmark. Unit tests and the standalone synthetic
demo remain explicit gates.

The demo does not require network access.

\section{17. Public Integration Contract}\label{17-public-integration-contract}

ARCANA exports model-bounded risk artifacts. Surrounding systems decide local actions.

Public integration rules:

\begin{itemize}
\tightlist
\item
  validate public schema before use;
\item
  validate model version, calibration profile, decision horizon, graph hash, evidence,
  verdict, and reason codes;
\item
  preserve original ARCANA reason codes;
\item
  preserve evidence source and hash binding;
\item
  keep public fields separate from integration-private fields;
\item
  reject or recompute stale, unsupported, mismatched, or uncertain artifacts;
\item
  keep A0 non-certifiable markers visible.
\end{itemize}

ARCANA remains enforcement-neutral. This paper does not define production enforcement,
private policy evaluation, private telemetry, or commercial certificate issuance.

ARCANA can be integrated with capability-bound execution systems, policy gateways, or
runtime brokers, but the public ARCANA reference track remains enforcement-neutral.

\section{18. Limitations}\label{18-limitations}

ARCANA is useful only inside its declared scope.

ARCANA is only as complete as the graph and evidence it is given.

Core limitations:

\begin{itemize}
\tightlist
\item
  the graph can omit real propagation paths;
\item
  calibration can be sparse, stale, wrong, or incomplete;
\item
  synthetic benchmark results do not prove production behavior;
\item
  A0 output is non-certifiable;
\item
  A1 output is not a commercial certificate basis in this public track;
\item
  A1 cannot support non-demo certificate-like artifacts in the public v0.2 contract;
\item
  A2/A3 language remains conditional on separate review, evidence coverage, freshness,
  horizon compatibility, non-synthetic evidence, expiry, and revocation handling;
\item
  synthetic \path{public_benchmark} evidence cannot qualify A2; only the distinct
  reviewed \path{empirical_public_benchmark} class can do so;
\item
  FastGate depends on valid assumptions and fallback paths;
\item
  evidence hashes bind artifacts but do not prove evidence quality;
\item
  local policy may be stricter than ARCANA\textquotesingle s bounded result;
\item
  a valid artifact is not a substitute for operational review;
\item
  public examples do not represent private deployments.
\end{itemize}

No public release should proceed when claim boundaries are unclear.

\section{19. Related Work}\label{19-related-work}

ARCANA is positioned as a bounded autonomy-accounting layer, not as a replacement for
tool protocols, threat taxonomies, governance frameworks, runtime verification,
proof-carrying authorization, network propagation models, or financial tail-risk
methods.

{\def\LTcaptype{none} % do not increment counter
\begin{longtable}[]{@{}>{\raggedright\arraybackslash}p{0.31\linewidth}>{\raggedright\arraybackslash}p{0.62\linewidth}@{}}
\toprule\noalign{}
Area & Relationship to ARCANA \\
\midrule\noalign{}
\endhead
\bottomrule\noalign{}
\endlastfoot
Secure message exchange & AXCP defines a separate protocol-oriented security layer for
agent communication {[}1{]}. ARCANA does not replace message security and does not
require AXCP; it accounts for model-bounded autonomy at the operation and
capability-graph layer. \\
Model Context Protocol (MCP) & MCP standardizes client-server access to resources,
prompts, tools, and related protocol primitives {[}2{]}. ARCANA does not compete with
tool connectivity; it accounts for autonomy risk in operations that such connectivity
can enable. \\
OWASP Agentic AI threats & OWASP\textquotesingle s Agentic AI guidance frames emerging
agentic threats and mitigations through a threat-model lens {[}3{]}. ARCANA can
complement such taxonomies with model-bounded risk artifacts, reason codes, and
calibration-scoped decisions. \\
NIST AI Risk Management Framework & NIST AI RMF is a voluntary framework for managing AI
risks to individuals, organizations, and society {[}4{]}. ARCANA does not replace AI
RMF; it provides a narrower mechanism for bounded-under-model decisions in agentic
operations. \\
Nonnegative matrix theory & ARCANA uses standard spectral-radius and nonnegative-matrix
results {[}5, 8{]}, while binding their interpretation to an explicit graph, horizon,
calibration profile, and uncertainty policy. \\
Branching processes and network propagation & ARCANA borrows the subcritical intuition
of finite propagation pressure from branching-process theory {[}6{]}, without assuming
that real agentic systems satisfy a classical branching model. \\
VaR and expected shortfall & ARCANA separates propagation risk from loss-side tail
quantities informed by conditional Value-at-Risk methods {[}7{]}. \\
Runtime verification and authorization & Runtime verification monitors executions
against specified properties {[}9{]}. ARCANA can inform runtime checks, but the public
reference track does not define enforcement or authorization infrastructure. \\
\end{longtable}
}

\section{20. Related Public Artifacts}\label{20-related-public-artifacts}

This manuscript is grounded in the current public research/reference artifacts:

{\def\LTcaptype{none} % do not increment counter
\begin{longtable}[]{@{}>{\raggedright\arraybackslash}p{0.31\linewidth}>{\raggedright\arraybackslash}p{0.62\linewidth}@{}}
\toprule\noalign{}
Artifact & Role in this manuscript \\
\midrule\noalign{}
\endhead
\bottomrule\noalign{}
\endlastfoot
\path{docs/FORMAL_MODEL_v0.2.md} & Mathematical contract for graph, \texttt{K},
\texttt{L}, horizon, uncertainty, and decision rule. \\
\path{docs/CALIBRATION_METHODOLOGY_v0.1.md} & Calibration levels, evidence quality,
sparse evidence, and staleness. \\
\path{docs/REASON_CODES.md} & Public reason-code contract. \\
\path{schemas/} & Public machine-readable artifact contracts. \\
\path{examples/} & Synthetic examples and benchmark fixtures. \\
\path{docs/FASTGATE_DESIGN_v0.1.md} & Conservative FastGate semantics. \\
\path{docs/ARCANA_BENCH_v0.1.md} & Synthetic benchmark scenario suite and
metrics. \\
\path{docs/PUBLIC_INTEGRATION_CONTRACT_v0.1.md} & Enforcement-neutral import/export
guidance. \\
\path{docs/LIMITATIONS_v0.1.md} & Explicit public limitations and release gates. \\
\path{src/arcana/} & Public reference implementation. \\
\path{tests/} & Local reproducibility and contract tests. \\
\end{longtable}
}

\section{21. Reproducibility and
Falsifiability}\label{21-reproducibility-and-falsifiability}

The public repository makes the reference claims inspectable through four separate
gates:

\begin{enumerate}
\def\labelenumi{\arabic{enumi}.}
\tightlist
\item
  \texttt{make\ check} validates the public boundary, schemas, examples, and executable
  ARCANA-Bench contract.
\item
  \texttt{make\ test} exercises typed contracts, explicit reason-code branches,
  numerical edge cases, artifact semantics, benchmark tamper detection, and
  claim-language guardrails.
\item
  \texttt{make\ demo} produces the deterministic A0 non-certifiable example.
\item
  \texttt{make\ bench} executes all synthetic benchmark inputs through the reference
  evaluator and verifies the source-bound report hash.
\end{enumerate}

The benchmark report distinguishes measured, not-measured, and unavailable metrics. A
failed expected-versus-observed comparison, source-manifest mismatch, content-hash
mismatch, semantic artifact contradiction, or changed decision reason is therefore
visible rather than silently normalized.

These checks make the reference implementation reproducible and falsifiable at the
contract level. They do not constitute external academic peer review, empirical
deployment validation, or evidence that an unmodeled path is absent. Exact commit,
dependency, and artifact hashes belong in the accompanying release record so that this
manuscript does not make mutable operational claims.

\section{22. Conclusion}\label{22-conclusion}

ARCANA defines a model-bounded autonomy-accounting layer for agentic operations. Its
central discipline is to bind every bounded result to a declared model, calibration
profile, horizon, graph, uncertainty interval, evidence assumption, budget, and
reason-code contract while keeping propagation pressure separate from loss.

The public reference implementation demonstrates these contracts on deterministic
synthetic inputs, including executable negative controls and a conservative FastGate
prototype. It establishes reference-evaluator behavior, not real-world safety, empirical
calibration quality, production enforcement, or commercial certificate authority. Those
stronger claims require evidence and governance outside the public reference package.

Any LaTeX, PDF, archive, or venue-specific derivative of this manuscript must preserve
this claim boundary and the public/private product separation.

\clearpage
\section{23. References}\label{23-references}

\begin{enumerate}
\def\labelenumi{\arabic{enumi}.}
\tightlist
\item
  J. Elizondo Rodriguez, "AXCP: Adaptive eXchange Context Protocol - Secure
  Communication for Autonomous AI Agents," version 1.0, Zenodo, 2026.
  \url{https://doi.org/10.5281/zenodo.18475648}
\item
  Model Context Protocol Contributors, "Architecture overview," Model Context Protocol
  documentation. Accessed 2026-07-13.
  \url{https://modelcontextprotocol.io/docs/learn/architecture}
\item
  OWASP Agentic Security Initiative, "Agentic AI - Threats and Mitigations," OWASP GenAI
  Security Project. Accessed 2026-07-13.
  \url{https://genai.owasp.org/resource/agentic-ai-threats-and-mitigations/}
\item
  E. Tabassi, "Artificial Intelligence Risk Management Framework (AI RMF 1.0)," NIST AI
  100-1, National Institute of Standards and Technology, 2023.
  \url{https://doi.org/10.6028/NIST.AI.100-1}
\item
  E. Seneta, \emph{Non-negative Matrices and Markov Chains}, 2nd ed., Springer, 1981.
  \url{https://doi.org/10.1007/0-387-32792-4}
\item
  T. E. Harris, \emph{The Theory of Branching Processes}, Springer, 1963. ISBN
  978-3-642-51868-3.
\item
  R. T. Rockafellar and S. Uryasev, "Optimization of Conditional Value-at-Risk,"
  \emph{The Journal of Risk}, vol. 2, no. 3, pp. 21-41, 2000.
  \url{https://doi.org/10.21314/JOR.2000.038}
\item
  R. A. Horn and C. R. Johnson, \emph{Matrix Analysis}, 2nd ed., Cambridge University
  Press, 2012. \url{https://doi.org/10.1017/CBO9780511810817}
\item
  M. Leucker and C. Schallhart, "A Brief Account of Runtime Verification," \emph{The
  Journal of Logic and Algebraic Programming}, vol. 78, no. 5, pp. 293-303, 2009.
  \url{https://doi.org/10.1016/j.jlap.2008.08.004}
\end{enumerate}

\end{document}
